\section{Real-World Policy Ambiguity Prevalence}
\label{app:prevalence}

To ground our taxonomy in real-world policy documents, we conduct a prevalence audit measuring how frequently each ambiguity type occurs in operational policies from major companies.

\subsection{Data Collection}

We collected 24 publicly available policy documents from 9 companies spanning two domains: retail (5 companies: Target, IKEA, Nordstrom, Zappos, Apple; 11 documents) and airlines (4 companies: Delta, Emirates, Frontier, Southwest; 13 documents). Documents include return policies, shipping policies, terms of service, contracts of carriage, and customer service plans. All documents were retrieved from official company websites on May 23, 2026.

\paragraph{Selection bias.}
We attempted to collect policies from 19 companies total. Ten companies blocked automated retrieval: Amazon (HTTP 503), Walmart (CAPTCHA), Costco (timeout), Best Buy (timeout), REI (timeout), Home Depot (HTTP 403), Macy's (HTTP 403), Sephora (HTTP 403), Nike (HTTP 403), and Lowe's (HTTP 403).
The direction of selection bias is uncertain: larger companies may maintain more professionally drafted policies (reducing ambiguity) or more complex policy structures (increasing it). Our sample is therefore a convenience sample of accessible policies, not a representative sample of all corporate policies.

\subsection{Clause Extraction}

From each document, we extracted all \emph{operational clauses}---sentences or multi-sentence blocks that prescribe, constrain, or condition a specific agent action or decision. This yielded \textbf{407 clauses} across the 24 documents (mean 17.0 per document; range 5--75). Each clause was assigned a unique identifier for annotation.

\subsection{Dual-Judge Annotation}

Each clause was independently annotated by two LLM judges: Claude Sonnet 4.6 (Anthropic) and GPT-5.4 (OpenAI). For each clause, judges determined: (1)~whether the clause is ambiguous in a way that could cause an LLM agent to misinterpret the intended policy, and (2)~which of our six ambiguity types best describes the ambiguity (scopal, lexical, coreferential, incompleteness, authorization scope, or conditional precedence).

\paragraph{Inter-annotator agreement.}
Cohen's $\kappa$ for the binary ambiguous/not-ambiguous decision was \textbf{0.295} (fair agreement). Claude flagged 83.8\% of clauses as ambiguous versus GPT's 58.5\%, reflecting Claude's more conservative threshold for ``unambiguous.'' Among the 226 clauses both judges marked as ambiguous, type-level agreement was $\kappa = 0.382$ with 58.4\% exact match.

\subsection{Results}

\paragraph{Document-level prevalence.}
All 24/24 documents (100\%) contained at least one clause flagged as ambiguous by either judge; all 24/24 also contained at least one clause flagged by \emph{both} judges.

\paragraph{Clause-level prevalence.}

\begin{table}[h]
\centering
\small
\begin{tabular}{lcc}
\toprule
\textbf{Metric} & \textbf{Union} & \textbf{Intersection} \\
\midrule
Raw (ambiguous / total) & 353/407 (86.7\%) & 226/407 (55.5\%) \\
Document-weighted mean  & 89.1\%  & 58.7\% \\
\bottomrule
\end{tabular}
\caption{Clause-level ambiguity prevalence. Union: flagged by either judge. Intersection: flagged by both.}
\label{tab:clause_prevalence}
\end{table}

\paragraph{Company-level analysis.}

\begin{table}[h]
\centering
\small
\begin{tabular}{llrrr}
\toprule
\textbf{Company} & \textbf{Domain} & \textbf{Clauses} & \textbf{Ambig.\ (union)} & \textbf{Rate} \\
\midrule
Apple & Retail & 43 & 39 & 90.7\% \\
Delta & Airline & 65 & 57 & 87.7\% \\
Emirates & Airline & 19 & 19 & 100.0\% \\
Frontier & Airline & 85 & 73 & 85.9\% \\
IKEA & Retail & 71 & 60 & 84.5\% \\
Nordstrom & Retail & 11 & 7 & 63.6\% \\
Southwest & Airline & 81 & 68 & 84.0\% \\
Target & Retail & 12 & 11 & 91.7\% \\
Zappos & Retail & 20 & 19 & 95.0\% \\
\midrule
\multicolumn{3}{l}{Company-level mean (N=9)} & \multicolumn{2}{c}{87.0\%, 95\% CI [80.0\%, 92.5\%]} \\
\bottomrule
\end{tabular}
\caption{Company-level ambiguity rates. Bootstrap 95\% CI uses company (N=9) as the clustering unit (B=10{,}000).}
\label{tab:company_prevalence}
\end{table}

\paragraph{Type distribution.}

\begin{table}[h]
\centering
\small
\begin{tabular}{lrr}
\toprule
\textbf{Ambiguity Type} & \textbf{Count} & \textbf{\% of Ambiguous} \\
\midrule
Incompleteness          & 221 & 62.6\% \\
Scopal                  & 47 & 13.3\% \\
Lexical                 & 36 & 10.2\% \\
Conditional precedence  & 33 & 9.3\% \\
Coreferential           & 9 & 2.5\% \\
Authorization scope     & 7 & 2.0\% \\
\bottomrule
\end{tabular}
\caption{Distribution of ambiguity types across all flagged clauses (union criterion).}
\label{tab:type_distribution}
\end{table}

\subsection{Limitations of This Pilot}

\paragraph{LLM-generated annotations.}
Ambiguity annotations are LLM-generated ($\kappa = 0.295$); human validation of a stratified 50-clause subset is planned for the camera-ready version. Neither judge's labels should be treated as ground truth. The moderate inter-judge disagreement (Claude flags 83.8\% vs.\ GPT 58.5\%) likely reflects differing ambiguity thresholds rather than random noise, as type agreement among agreed-ambiguous clauses is substantially higher (58.4\% exact match).

\paragraph{Sample coverage.}
Our corpus covers 9 of 19 targeted companies; 10 companies blocked automated retrieval. Generalization to other languages, jurisdictions, or industry domains requires further study.
